\chapter{Installation}
\label{ch:installation}

This chapter installs ProPrep on your computer. It assumes you have never used
the Terminal before. Every command is shown exactly; type it, or copy and paste
it, and press \key{Return} (Mac) or \key{Enter} (Windows and Linux) at the end of
each line.

There are two ways to install ProPrep:
\begin{description}
  \item[Option A, the self-contained installer (recommended).] One file that
        contains ProPrep, AmberTools, MODELLER and everything else. No other
        software is needed. Use this if you are not sure.
  \item[Option B, the conda installer script.] For people who already use
        \term{conda}, a program that installs scientific software into separate
        folders called \term{conda environments}. It downloads everything
        during the installation.
\end{description}

Both options give you the same ProPrep. After installing, add your MODELLER key
(Section~\ref{sec:installation-modeller}) and run the check in
Section~\ref{sec:installation-check}.

\section{Option A: the self-contained installer}
\label{sec:installation-optiona}

\begin{behindscenes}
The computer used for this manual already had conda, so it was set up with
Option~B. The Option~A steps below follow ProPrep's official installation
guide (\file{INSTALL.md}); the messages quoted are the ones that guide lists.
\end{behindscenes}

You will download one file of about 1 to 1.5\,GB. You need about 8\,GB of free
disk space and 10 to 15 minutes. Once the file is downloaded, the installation
itself needs no internet connection. You end up with a folder called
\file{ProPrep} in your home folder.

\subsection{Choose and download the installer}

\begin{tabularx}{\linewidth}{@{}lL@{}}
\toprule
\textbf{Your computer} & \textbf{Installer file} \\
\midrule
Mac with an Apple chip (M1, M2, M3, \ldots) & \file{ProPrep-<version>-MacOSX-arm64.sh} \\
Mac with an Intel chip & \file{ProPrep-<version>-MacOSX-x86_64.sh} \\
Windows 10 or 11 (through WSL) & \file{ProPrep-<version>-Linux-x86_64.sh} \\
Linux & \file{ProPrep-<version>-Linux-x86_64.sh} \\
\bottomrule
\end{tabularx}

\medskip
\texttt{<version>} stands for the version number, for example
\texttt{1.23.0}.

\begin{tip}
Not sure which Mac you have? Click the Apple menu (top-left corner) and choose
\textsf{About This Mac}. \textsf{Chip: Apple M\ldots} means an Apple chip;
\textsf{Processor: Intel\ldots} means an Intel chip.
\end{tip}

\begin{steps}
  \item Go to \url{https://github.com/mjgplab/proprep/releases/latest}.
  \item Scroll down to \textsf{Assets} and click the file name for your
        computer. Your browser saves it in your \file{Downloads} folder. The
        file is large, so allow a few minutes.
\end{steps}

\subsection{On a Mac}
\label{sec:installation-mac}

\begin{steps}
  \item Open the Terminal: press \key{Command}+\key{Space}, type
        \typed{Terminal}, and press \key{Return}. A window with a text prompt
        appears.
  \item Type the line for your Mac. Apple chip:
\begin{terminal}
bash ~/Downloads/ProPrep-<version>-MacOSX-arm64.sh -b -p ~/ProPrep
\end{terminal}
        Intel chip:
\begin{terminal}
bash ~/Downloads/ProPrep-<version>-MacOSX-x86_64.sh -b -p ~/ProPrep
\end{terminal}
        The sign \texttt{\textasciitilde} means your home folder. The option
        \texttt{-b} means ``accept the license and ask no questions'';
        \texttt{-p \textasciitilde/ProPrep} is where ProPrep is installed.
  \item Wait 2 to 5 minutes. The installer prints \prompt{Unpacking
        payload...} and a list of packages, then \prompt{ProPrep is
        installed.}
  \item If macOS asks whether Terminal may access your Downloads folder,
        click \textsf{OK}. The installer only reads the file you downloaded.
\end{steps}

\begin{tip}
You do not have to type the long file name. Type
\texttt{bash \textasciitilde/Downloads/ProPrep-} and press the \key{Tab} key;
the Terminal completes the name for you.
\end{tip}

\subsection{On Windows (through WSL)}
\label{sec:installation-windows}

ProPrep runs on Linux and macOS. On Windows it runs inside the Windows
Subsystem for Linux (WSL), a feature built into Windows that runs a small Linux
system called \term{Ubuntu} next to Windows. You set up WSL once. After that,
ProPrep opens in your normal Windows web browser.

\paragraph{Set up WSL (once).}
\begin{steps}
  \item Click \textsf{Start}, type \typed{PowerShell}, right-click
        \textsf{Windows PowerShell} and choose \textsf{Run as administrator}.
        Click \textsf{Yes} if Windows asks for permission.
  \item In the blue window, type:
\begin{terminal}
wsl --install
\end{terminal}
        Windows downloads and installs WSL and Ubuntu. This takes several
        minutes.
  \item Restart your computer when Windows asks you to.
  \item After the restart, an Ubuntu window opens and says
        \prompt{Installing, this may take a few minutes...}. It then asks you
        to choose a user name (lowercase, no spaces) and a password. The
        password does not show while you type; that is normal. Remember it.
        If no Ubuntu window opens, click \textsf{Start}, type
        \typed{Ubuntu}, and open it.
\end{steps}

You now have an Ubuntu window with a prompt ending in \texttt{\$}. This is
your Terminal on Windows. Type every command below in it.

\paragraph{Install ProPrep.} Your Windows \file{Downloads} folder is visible
from Ubuntu as\newline \file{/mnt/c/Users/<YourWindowsName>/Downloads}.\newline Type this,
replacing \texttt{<YourWindowsName>} with your Windows user name and
\texttt{<version>} with the version number:

\begin{terminal}
bash /mnt/c/Users/<YourWindowsName>/Downloads/ProPrep-<version>-Linux-x86_64.sh \
  -b -p ~/ProPrep
\end{terminal}

The backslash at the end of the first line means ``the command continues on
the next line''. You can also type it all on one line without the backslash.
If you do not know your Windows user name, type \texttt{ls /mnt/c/Users} and
press \key{Enter} to see the list. The installer takes 2 to 5 minutes and ends
with \prompt{ProPrep is installed.}

\paragraph{Where are my files?} ProPrep saves projects in your Ubuntu home
folder. To see them in Windows Explorer, type this in the Explorer address bar
(with your Ubuntu user name): \verb|\\wsl$\Ubuntu\home\<your-ubuntu-username>|

\begin{watchout}
WSL needs Windows 10 version 2004 or later, or Windows 11, and about 10\,GB of
free disk space in total. If \texttt{wsl --install} reports an error, the usual
causes are an old Windows 10 (run Windows Update) or virtualization switched
off in the computer's firmware; your IT support can switch it on.
\end{watchout}

\subsection{On Linux}
\label{sec:installation-linux}

Open a terminal and type:

\begin{terminal}
bash ~/Downloads/ProPrep-<version>-Linux-x86_64.sh -b -p ~/ProPrep
\end{terminal}

The installer takes 2 to 5 minutes and puts ProPrep in \file{~/ProPrep}. It is
built for 64-bit Intel and AMD processors (x86-64) and has been tested on Rocky
Linux.

\begin{tip}
On a computer without a screen that you reach over \texttt{ssh} (for example a
cluster), \texttt{proprep-web} prints the exact \texttt{ssh -L} command to run
on your own computer. After that, \url{http://localhost:8000/} opens ProPrep in
your local browser.
\end{tip}

\subsection{Starting ProPrep after Option A}

\begin{terminal}
~/ProPrep/bin/proprep-web
\end{terminal}

Your web browser opens with ProPrep inside it. On Windows, open
\url{http://127.0.0.1:8000/} in your Windows browser yourself; the Ubuntu
window prints a line \prompt{ProPrep Web Shell listening on
http://127.0.0.1:8000/}. Leave the Terminal window open while you work. To use
ProPrep in the Terminal instead of the browser, type
\texttt{\textasciitilde/ProPrep/bin/proprep}. Chapter~\ref{ch:starting}
explains both.

\section{Option B: the conda installer script}
\label{sec:installation-optionb}

Use this option if you already have conda. If you do not, install Miniforge
first from\newline \url{https://github.com/conda-forge/miniforge#download},\newline close the
Terminal and open a new one. Without conda the script stops with
\prompt{ERROR: conda is not installed or not in your PATH.}

The script creates a conda environment named \texttt{ProPrep} with Python 3.12,
installs ProPrep 1.23.0 and its dependencies, and checks the result.

\begin{steps}
  \item Download the script and run it:
\begin{terminal}
curl -fsSLO \
  https://raw.githubusercontent.com/mjgplab/proprep/main/install_proprep.sh
bash install_proprep.sh
\end{terminal}
  \item The script asks for your MODELLER key. The characters you type show
        as \texttt{*}. If you have no key yet, press \key{Enter}; you can add
        it later (Section~\ref{sec:installation-modeller}).
  \item Wait. The script downloads several gigabytes, so it takes several
        minutes. It prints the package list as it goes.
  \item At the end it checks the version and prints how to start ProPrep.
\end{steps}

This is what we saw on our test Mac (long parts left out):

\begin{screen}
============================================================
  ProPrep Installer
============================================================

MODELLER requires a license key (free for academic use).
Register at: https://salilab.org/modeller/registration.html

Enter your MODELLER license key (or press Enter to skip):

NOTE: Skipping MODELLER key. Homology modeling will not be
available until you configure the key manually. See the
install guide for instructions.
...
Verifying installation...
  ProPrep installed successfully: ProPrep v1.23.0 - Proper Protein Preparation...

============================================================
  Installation complete!
============================================================

  To use ProPrep:

    conda activate ProPrep
    source $CONDA_PREFIX/amber.sh
    proprep              # interactive CLI
    proprep-web          # browser-based UI
\end{screen}

``Homology modeling'' in the note means all the tools that need MODELLER,
including structure repair. The environment took 4.8\,GB on our Mac.

\paragraph{Starting ProPrep after Option B.} Every time you open a new
Terminal window, type these two lines first, then \texttt{proprep-web} or
\texttt{proprep}:

\begin{terminal}
conda activate ProPrep
source $CONDA_PREFIX/amber.sh
\end{terminal}

The first line switches to ProPrep's environment. The second makes sure the
Amber programs come from that environment
(Section~\ref{sec:installation-amberhome}).

\begin{tip}
The README also offers a one-line version that downloads and runs the script
in one go:
\texttt{curl -fsSL https://raw.githubusercontent.com/mjgplab/proprep/main/install\_proprep.sh | bash}.
It does the same thing. Downloading first lets you look at the script before you
run it.
\end{tip}

\section{Adding your MODELLER license key}
\label{sec:installation-modeller}

ProPrep uses \term{MODELLER} to build missing parts of a protein, such as
loops that were not seen in the experiment. MODELLER is installed with ProPrep,
but it only works with a personal \term{license key}, a code that proves you
registered. The key is free for academic use.

\begin{steps}
  \item Register at\newline \url{https://salilab.org/modeller/registration.html}\newline
        (name, email, institution; choose the academic license). The key
        arrives by email within a few minutes.
  \item In the Terminal, type this line, replacing \texttt{YOUR\_KEY} with
        your key and keeping the quotation marks:
\begin{terminal}
mkdir -p ~/.proprep && echo 'YOUR_KEY' > ~/.proprep/modeller_key
\end{terminal}
\end{steps}

ProPrep reads the key from the file \file{~/.proprep/modeller_key} every time
it starts. If you mistype it, run the same line again with the correct key.
The key stays on your computer. As an alternative, you can set the key in the
Terminal with \texttt{export KEY\_MODELLER='YOUR\_KEY'} before starting
ProPrep; this setting wins over the file.

\begin{behindscenes}
We had no MODELLER key while writing this manual, so the steps above come from
ProPrep's installation guide and source code. We could not show the screens of
a working MODELLER repair.
\end{behindscenes}

\subsection{The ``MODELLER not found'' message}
\label{sec:installation-modeller-msg}

Without a working key, every start of ProPrep begins with this message:

\begin{screen}
MODELLER not found. Structure repair (missing loops/residues) will be unavailable.
  To install: conda install -c salilab modeller
  A free academic license is required: https://salilab.org/modeller/registration.html
  After registering, set your license key:
    export KEY_MODELLER='your_license_key_here'
  Or edit the config.py file shown during conda installation.
Successfully imported core modules
\end{screen}

\begin{watchout}
In version 1.23.0 this message says ``not found'' and ``To install'' even
when MODELLER \emph{is} installed and only the key is missing. That was the
case on our test computer. Do not install MODELLER again. Add the key as shown
above. Everything else in ProPrep works without it; you can ignore the message
if you never need to repair loops or missing residues.
\end{watchout}

\section{Checking that it works}
\label{sec:installation-check}

With Option A, type:

\begin{terminal}
~/ProPrep/bin/proprep --version
\end{terminal}

With Option B, type \texttt{conda activate ProPrep} first, then
\texttt{proprep --version}. You should see:

\begin{screen}
ProPrep v1.23.0 - Proper Protein Preparation workflow manager
\end{screen}

Then start ProPrep (\texttt{proprep-web} or \texttt{proprep}). If the
``MODELLER not found'' message does not appear at the top, your key was found.
To stop ProPrep, close the browser tab (ProPrep shuts down by itself a moment
later) or press \key{Control}+\key{C} in the Terminal.

\section{Making sure ProPrep uses its own Amber}
\label{sec:installation-amberhome}

ProPrep runs Amber programs such as tLEaP and sander. To find them, it looks at
a setting called \term{AMBERHOME}, which names the folder of an Amber
installation. If AMBERHOME is already set when ProPrep starts, ProPrep keeps
it. So if another Amber or AmberTools installation is active in your Terminal,
ProPrep uses that one instead of its own copy, and it does not tell you.

This happened on our test computer: an older conda environment with
AmberTools~25 was active, and tLEaP, sander and the force-field files all came
from there. You can spot it in tLEaP's messages, which show the program's
folder:

\begin{screen}
Warning 1:
  /opt/homebrew/Caskroom/miniforge/base/envs/etanalyze_env/bin/teLeap: Warning!
  Close contact of 1.306 angstroms between nonbonded atoms H and H1
\end{screen}

Here \file{etanalyze_env} is the other environment, not \file{ProPrep}.

\paragraph{Check before you start.} In the Terminal window you will start
ProPrep from, type:

\begin{terminal}
echo $AMBERHOME
\end{terminal}

\begin{itemize}
  \item An empty line, or a folder inside \file{ProPrep} (for example
        \file{~/ProPrep} or \file{.../envs/ProPrep}): all is well.
  \item Any other folder: another Amber is active. Fix it in this window, then
        start ProPrep from the same window. With Option A, type:
\begin{terminal}
unset AMBERHOME
\end{terminal}
        With Option B, type:
\begin{terminal}
conda activate ProPrep
source $CONDA_PREFIX/amber.sh
\end{terminal}
        Then type \texttt{echo \$AMBERHOME} again to check.
\end{itemize}

\begin{tip}
If every new Terminal window shows the other folder, a start-up file (such as
\file{~/.zshrc} or \file{~/.bashrc}) sets it. You can keep it, and run the fix
above each time, or ask whoever set up the other Amber to help you change it.
\end{tip}

\begin{behindscenes}
To make sure our results did not depend on the other installation, we ran
ProPrep's ubiquitin tLEaP input again with only ProPrep's own AmberTools~26
active. tLEaP finished with \prompt{Errors = 0; Warnings = 13} and wrote a
complete topology, read from ProPrep's own force-field files.
\end{behindscenes}

\section{Updating ProPrep}
\label{sec:installation-update}

ProPrep checks GitHub once a day, when it starts, for a newer release. It sends
nothing about you. If there is one, it prints a line such as
\prompt{ProPrep X.Y.Z is available (you have 1.23.0).} with the update command
for your kind of installation. (We did not see this line, because 1.23.0 was
the newest release.) You can switch the check off in Preferences
(Chapter~\ref{ch:menus}).

\paragraph{Option A.} Download the new installer and run it as before, with
\texttt{-u} added. It updates the \file{~/ProPrep} folder in place:

\begin{terminal}
bash ~/Downloads/ProPrep-<version>-<platform>.sh -b -u -p ~/ProPrep
\end{terminal}

\paragraph{Option B.} Run \texttt{bash install\_proprep.sh} again (download
the newest copy first). It finds your environment and asks:

\begin{screen}
Environment 'ProPrep' already exists.

  1. Update ProPrep (keep existing environment)
  2. Full reinstall (remove and recreate environment)
  3. Cancel

Choose [1/2/3] (1):
\end{screen}

Choose \key{1} to update. Choose \key{2} if the environment seems damaged.

Your projects, your MODELLER key and your settings live outside the
installation, so updating does not touch them.

\section{Removing ProPrep}
\label{sec:installation-remove}

\begin{itemize}
  \item Option A: delete the folder \file{ProPrep} in your home folder, or type
        \texttt{rm -rf \textasciitilde/ProPrep} in the Terminal.
  \item Option B: type \texttt{conda env remove -n ProPrep}.
\end{itemize}

Your projects stay where you saved them. Your MODELLER key, settings and
personal parameter library stay in the hidden folder \file{~/.proprep}; delete
it too if you want to remove everything.

\section{If something goes wrong}
\label{sec:installation-problems}

\begin{longtable}{@{}>{\raggedright\arraybackslash}p{0.36\linewidth}>{\raggedright\arraybackslash}p{0.58\linewidth}@{}}
\toprule
\textbf{What you see} & \textbf{What to do} \\
\midrule
\endhead
\leavevmode\prompt{No such file or directory} when running the installer &
The file name or folder is wrong, or the download did not finish. Use the
\key{Tab} key to complete the name. The file should be about 1\,GB or more. \\
\leavevmode\prompt{Permission denied} &
The line must begin with \texttt{bash}. Retype it. \\
\leavevmode\prompt{ERROR: File or directory already exists} &
ProPrep is already installed. To update it, add \texttt{-u}
(Section~\ref{sec:installation-update}). \\
\leavevmode\prompt{No space left on device} &
Free at least 8\,GB and run the installer again. \\
\texttt{wsl --install} needs a restart, then nothing happens &
Click \textsf{Start}, type \typed{Ubuntu}, and open it. The set-up continues
there. \\
Browser says \prompt{This site can't be reached} at 127.0.0.1:8000 &
ProPrep is not running. Start \texttt{proprep-web} again. If port 8000 is
busy, ProPrep picks the next one, for example
\leavevmode\prompt{(port 8420 was busy; using 8421)}. Use the address it prints. \\
Antivirus or company software blocks the installer &
The \file{.sh} file is a plain script that unpacks files into your home
folder. Ask your IT support to allow it. \\
\leavevmode\prompt{ERROR: conda is not installed or not in your PATH.} (Option B) &
Install Miniforge, close the Terminal, open a new one, and run the script
again. \\
\leavevmode\prompt{ERROR: expected ProPrep 1.23.0 but conda resolved:} (Option B) &
Conda chose a different version. Read the conflict messages above it. On a
Mac, the script suggests running it again as
\texttt{SYSTEM\_VERSION\_COMPAT=0 bash install\_proprep.sh}. (Not seen in our
tests.) \\
\leavevmode\prompt{MODELLER not found} at every start &
Usually only the license key is missing. Add it
(Section~\ref{sec:installation-modeller}). \\
\texttt{proprep: command not found} (Option B) &
The ProPrep environment is not active. Type \texttt{conda activate ProPrep}
first. \\
Results or messages mention a different Amber folder &
Another Amber installation is active. See
Section~\ref{sec:installation-amberhome}. \\
\texttt{proprep-library} stops with \prompt{bad interpreter} (Option B) &
In version 1.23.0 this helper points to the wrong Python. Use
\texttt{python -m proprep.utils.user\_library\_snapshots} instead, with the same
options (Chapter~\ref{ch:forcefields}). \\
\bottomrule
\end{longtable}

For other problems, see Chapter~\ref{ch:troubleshooting}.
